%==============================================================================
%  EHRJEPA v2 -- draft (Oct 2026). Build: pdflatex main; bibtex main; pdflatex x2
%  Red [TODO: ...] markers = results of the state-model gate, not yet available.
%==============================================================================
\documentclass[10pt,letterpaper]{article}
\usepackage[margin=1in]{geometry}
\usepackage[T1]{fontenc}
\usepackage[utf8]{inputenc}
\usepackage{lmodern}
\usepackage{microtype}
\usepackage{amsmath,amssymb}
\usepackage{booktabs}
\usepackage{multirow}
\usepackage{xcolor}
\usepackage{caption}
\usepackage[numbers,sort&compress]{natbib}
\usepackage[hidelinks]{hyperref}
\usepackage{url}
\usepackage{enumitem}
\emergencystretch=3em
\captionsetup{font=small,labelfont=bf}

\newcommand{\todo}[1]{\textcolor{red}{\textbf{[TODO: #1]}}}
\newcommand{\code}[1]{\texttt{\small #1}}
\newcommand{\ci}[2]{\,[{#1},\,{#2}]}

\title{\textbf{Latent Prediction Helps EHR Foundation Models Only Next to a Token Loss,\\
and Only Fine-Tuning Shows It:\\
A Careful Public-Data Study and an Action-Conditioned Patient-State Model}}
\author{Abel Yagubyan\\
\small University of California, Berkeley\\
\small \texttt{abelyagubyan@berkeley.edu}}
\date{Draft, October 2026}

\begin{document}
\maketitle

%==============================================================================
\begin{abstract}
Joint-embedding predictive architectures (JEPAs) have been proposed as
pretraining objectives for structured electronic health records (EHRs), alone
and as an auxiliary term to next-token prediction. Published EHR evidence is
mixed and rests almost entirely on frozen linear probes, single seeds and
private or text-serialized data. We train small (10M-parameter) transformers
from scratch on coded MIMIC-IV events and compare objectives under three
seeds, a leakage-checked count-feature gradient-boosting baseline, paired
subject-level bootstrap intervals, and both frozen probing and end-to-end
fine-tuning. Three findings. (i) Every objective without a dense next-code loss
trails next-code autoregression (AR) by 0.06--0.09 mean AUROC, agreeing with
recent reports that latent prediction alone underperforms token-level
objectives. (ii) Adding a multi-horizon next-latent term to AR (a hybrid)
improves fine-tuned performance over AR by $+0.008$ AUROC
\ci{+0.002}{+0.014}; against gradient boosting the hybrid gains $+0.013$
\ci{-0.001}{+0.028}, which is not significant. (iii) The frozen probe ranks the
hybrid \emph{below} gradient boosting (0.803 vs 0.808) while fine-tuning ranks it
above (0.822), so the protocol used by prior EHR-JEPA work can hide the effect
it is meant to measure. Following the suggestion that EHR latent targets
should be patient states whose transitions are shaped by care decisions, and
prior action-conditioned EHR JEPA work on text encoders, we test a small
code-token action-conditioned patient-state objective with a treatment-blind control: 4-hour state
representations predicted $k\in\{1,2,4\}$ windows ahead from the current state and
the intervening treatments, trained jointly with the next-code loss.
\todo{gate result: state\_ac vs hybrid / AR / MLM / GBM, fine-tuned mean-of-6 with CI}
\todo{action-sensitivity result: state\_ac vs treatment-blind control}
Code and configurations are public.
\end{abstract}

%==============================================================================
\section{Introduction}
\label{sec:intro}

Foundation models for structured EHRs are mostly trained by predicting the
next medical code \citep{clmbr,ethos,cehrgpt,delphi2m,motor} or masked codes
\citep{behrt,medbert}. JEPAs \citep{ijepa,vjepa2,lejepa} instead predict the
\emph{representation} of withheld content, which in vision and video gives
representations that are cheaper to learn and more abstract than pixel
reconstruction, and V-JEPA~2-AC \citep{vjepa2} uses an action-conditioned
latent predictor as a world model for planning. Whether this carries over to
EHRs is unresolved. SMB-Structure \citep{smbstructure} combined next-token
prediction with an EMA-target JEPA term on a text-serialized 8B language model
and reported gains of up to about $+0.013$ on private data under linear
probing. Clin-JEPA \citep{clinjepa} reported gains over LightGBM with a
multi-phase co-training recipe on an LLM text encoder and already trains an
action-conditioned latent world model on hourly ICU state--action pairs. Gupta et
al.\ \citep{imperialjepa} found that three JEPA formulations on MIMIC-IV train
stably but consistently underperform masked language modelling (MLM) under
linear probing, and proposed that EHR latent targets be aggregated
\emph{patient states} whose transitions are modelled, noting that future
observations are partly determined by clinical decisions rather than by the
patient alone.

Two methodological gaps make these results hard to combine. First, their downstream comparisons rest on frozen
representations with a linear probe; none compares the probe ranking of objectives with
the ranking after end-to-end fine-tuning, where downstream use actually happens. Second, none reports
seeds, paired intervals, or a tuned tabular baseline audited for label
leakage, so differences of $\pm 0.01$ AUROC---the size of every reported gain---cannot be
distinguished from noise.

\paragraph{Contributions.}
\begin{enumerate}[leftmargin=*,itemsep=2pt,topsep=2pt]
\item \textbf{A careful, public-data comparison.} MIMIC-IV \citep{mimiciv} in
MEDS format \citep{meds}, one token per coded event, 10M-parameter models
trained from scratch at matched token budgets, three seeds, paired
subject-level bootstrap, a count-feature XGBoost baseline checked for label
leakage, and both frozen-probe and fine-tuned evaluation on identical anchors
(\S\ref{sec:eval}).
\item \textbf{Latent prediction helps only as an auxiliary to a token loss.}
Latent-only objectives lose to AR by 0.06--0.09, agreeing with
\citet{imperialjepa}. A next-latent term \emph{added} to AR (the hybrid; the
same family of combination as \citep{smbstructure}) beats AR significantly under fine-tuning,
$+0.008$ \ci{+0.002}{+0.014} (\S\ref{sec:results}).
\item \textbf{The probe hides the gain.} Frozen probing ranks the hybrid below
tuned gradient boosting; fine-tuning reverses the order. The hybrid's margin
over gradient boosting, $+0.013$ \ci{-0.001}{+0.028}, is not significant on our
1{,}346-patient labelled test set, and we do not claim it.
\item \textbf{A code-token patient-state model with a treatment-blind control.}
Action-conditioned latent prediction on MIMIC-IV has been proposed before
with an 8B text encoder \citep{clinjepa}. We test a small, from-scratch,
code-token version of the patient-state direction proposed by
\citet{imperialjepa}: events pooled into 4-hour states, treatments separated
from observations and given to the predictor as inputs, trained jointly with
the next-code loss, and compared against an otherwise identical
treatment-blind control and an MLM baseline (\S\ref{sec:state}). The model
is predictive, not causal. \todo{one-line gate outcome}
\end{enumerate}
We also outline a public benchmark that checks action-conditioned EHR models
against effects known from randomized trials (\S\ref{sec:discussion}); it is
planned work and no results are reported for it.

%==============================================================================
\section{Related work}
\label{sec:related}

\paragraph{EHR foundation models.}
Next-code models on coded events include CLMBR \citep{clmbr}, the
time-to-event model MOTOR \citep{motor}, ETHOS \citep{ethos}, CEHR-GPT
\citep{cehrgpt}, EHRMamba \citep{ehrmamba} and Delphi-2M \citep{delphi2m};
BEHRT and Med-BERT use masked-code objectives \citep{behrt,medbert}. Scaling,
context length and tokenization have been studied
\citep{ehrscaling,contextclues,tokenization}. EHRSHOT \citep{ehrshot} and FoMoH
\citep{fomoh} standardize evaluation, and MEDS \citep{meds}, ACES \citep{aces}
and MEDS-Tab \citep{medstab} standardize data, cohorts and tabular baselines;
we use all three.

\paragraph{The JEPA family.}
I-JEPA \citep{ijepa} predicts masked image-block representations from an EMA
target encoder; V-JEPA~2 \citep{vjepa2} extends this to video and, with an
action-conditioned predictor (V-JEPA~2-AC), to planning. LeJEPA
\citep{lejepa} replaces architectural anti-collapse devices with the SIGReg
distribution penalty, which we use. LLM-JEPA \citep{llmjepa} adds a JEPA term
to a language model's next-token loss; SALT \citep{salt} uses frozen teachers;
LeNEPA \citep{lenepa} and TS-JEPA \citep{tsjepa} predict next latents in time
series.

\paragraph{JEPAs for EHRs.}
SMB-Structure \citep{smbstructure} is, to our knowledge, the first
next-token\,+\,JEPA hybrid for longitudinal EHRs, framed as a world-model
training paradigm; it serializes records as text for an 8B LLM fine-tuned with
LoRA, trains on private oncology data and INSPECT, evaluates by linear probe, and
reports that the joint (non-curriculum) hybrid can hurt. Clin-JEPA
\citep{clinjepa} co-trains a Qwen3-8B text encoder with a latent predictor on a
MIMIC-IV ICU cohort, and adds an action-conditioned latent world model
(hourly state--action text pairs, a $\sim$92M block-causal predictor,
evaluated by rollouts), without trial-based checks or a public benchmark. Gupta et al.\ \citep{imperialjepa} compare three JEPA
formulations with MLM on MIMIC-IV under linear probing and find JEPA
consistently worse. Our study differs in using coded events with a small model
trained from scratch on public data, in evaluating fine-tuned models with
seeds and paired intervals, and, for the state model, in using code tokens
with a treatment-blind control rather than a text encoder.

\paragraph{World models and treatment effects.}
Latent world models predict future latent state given actions
\citep{worldmodels,vjepa2}. Counterfactual outcome models over time, such as
CRN \citep{crn} and the Causal Transformer \citep{causaltransformer}, adjust
for time-varying confounding by treatment assignment. Target-trial emulation
\citep{hernan2016} and its large-scale replication against randomized trials
\citep{rctduplicate} provide the external reference we plan to use for
evaluating action-conditioned EHR models (\S\ref{sec:discussion}).

%==============================================================================
\section{Data}
\label{sec:data}

\paragraph{Source and format.}
MIMIC-IV v3.1 \citep{mimiciv} is lowered to MEDS with \code{meds\_etl}. Shards
carry the MEDS core columns (\code{subject\_id}, \code{time}, \code{code},
\code{numeric\_value}, \code{text\_value}). Splits are a deterministic hash of
the subject identifier (80/10/10 train/tuning/held-out).

\paragraph{Three pretraining pools, one test set.}
We build three subject-nested pools: an 8\% subset (every 12th patient,
30{,}386 subjects; 24{,}240 train subjects, 45.3M train events), a 25\% subset
(91{,}157 subjects; 72{,}900 train subjects, 138.0M events) and the full
database (291{,}495 train subjects, 546.4M events; 248{,}743 with at least 16
events are used for pretraining). The vocabulary and value quantizer are fit on
the 8\% train split and reused for the larger pools (99.91\% of full-pool
codes map directly), so pools differ only in subject count. All 3{,}038
held-out subjects of the 8\% split lie inside the larger pools' held-out
splits, and none is in any pretraining split. \emph{Every downstream number in
this paper is scored on these 8\% held-out subjects, with probes, fine-tuning
heads and baselines fit on the 8\% train labels}; this isolates the effect of
pretraining data from that of label count.

\paragraph{Tokenization.}
One event becomes one token. Codes with at least five training occurrences get
their own id; rarer codes roll up to the nearest frequent ancestor (ICD
prefixes, NDC labeller, etc.) before \code{[UNK]}, giving 18{,}736 ids. Each
token also carries a per-code decile bin and $z$-score of its numeric value,
age, and the log time since the previous event. These are summed into the
input embedding through Fourier-feature MLPs; position enters separately
through rotary embeddings. Models read windows of up to 512 events.

\paragraph{Actions.}
For the state model (\S\ref{sec:state}) we mark as \emph{actions} the events
whose codes denote treatments: prescriptions and administrations (NDC and
MIMIC drug codes), procedures (ICD-9 procedure, ICD-10-PCS, HCPCS) and ICU
input events, plus a curated list of procedure and ingredient item ids. All
other events are \emph{observations}. On the MIMIC-IV demo cache about 7.5\% of
events are actions.

%==============================================================================
\section{Method}
\label{sec:method}

\paragraph{Encoder.}
All models share a pre-norm transformer encoder $f_\theta$ (6 layers, width
256, 4 heads, SwiGLU feed-forward, rotary positions; $\approx$10.3M trainable
parameters), causal for AR, hybrid and state objectives and bidirectional for
masked objectives. Let $h_i$ be the encoder output at event $i$ and
$E\in\mathbb{R}^{V\times d}$ the code embedding table.

\paragraph{Next-code AR.}
A tied head scores the next code:
\begin{equation}
\mathcal{L}_{\mathrm{code}} = -\frac{1}{N}\sum_i
\log\frac{\exp(E_{c_{i+1}}^{\top}h_i)}{\sum_{c'}\exp(E_{c'}^{\top}h_i)} .
\end{equation}
Only the code is predicted; values and times are inputs.

\paragraph{Latent-only objectives.}
We compare four objectives without a dense next-code loss on the encoder's own states:
(a) masked-span JEPA with a bidirectional encoder and EMA target; (b) the same
architecture trained only to reconstruct masked codes from predicted latents;
(c) causal next-latent prediction (below, with $\lambda_{\mathrm{code}}=0$);
(d) prediction of mean target latents over 30- and 365-day windows.

\paragraph{Target encoder and SIGReg.}
Latent targets come from an EMA copy $f_{\bar\theta}$,
$\bar\theta \leftarrow m(t)\bar\theta + (1-m(t))\theta$, with $m$ rising
linearly from 0.996 to 1, under stop-gradient, and are normalized by a
parameter-free LayerNorm: $\tilde z_j = \mathrm{LN}(f_{\bar\theta}(x)_j)$.
Collapse is discouraged by SIGReg \citep{lejepa}: embeddings are projected onto
$K=256$ random directions and each projection's empirical characteristic function is
penalized for deviating from a standard normal's.

\paragraph{Hybrid.}
For horizons $\mathcal{H}=\{1,4,16\}$ events, a per-horizon MLP head $g^{(k)}$
receives $h_i$ plus a time embedding of the target event and predicts
$\tilde z_{i+k}$:
\begin{equation}
\mathcal{L}_{\mathrm{lat}}=\frac{1}{|\mathcal{H}|}\sum_{k\in\mathcal{H}}
\frac{1}{|\mathcal{V}_k|}\sum_{i\in\mathcal{V}_k}
\ell_\beta\big(g^{(k)}(h_i+u(t_{i+k})),\,\tilde z_{i+k}\big),
\qquad
\mathcal{L}_{\mathrm{hyb}} = \mathcal{L}_{\mathrm{lat}} + 0.1\,\mathcal{L}_{\mathrm{code}} + 0.05\,\mathrm{SIGReg},
\end{equation}
with $\ell_\beta$ the smooth-$L_1$ loss averaged over feature dimensions. Setting the latent
weight to zero recovers AR exactly (a unit test pins this), so the hybrid is a
controlled addition to AR.

\paragraph{MLM baseline.}
Bidirectional masked-code prediction with 15\% masking (80/10/10 replacement),
the baseline of \citet{imperialjepa}, on the same encoder.

\subsection{Action-conditioned patient-state objective}
\label{sec:state}

Following the proposal of \citet{imperialjepa}, the prediction unit is a
patient \emph{state}, not an event, and following their observation that future
observations are shaped by care decisions, treatments are given to the
predictor as inputs rather than left for it to infer.

\paragraph{States.}
Each 512-event crop is cut into calendar windows of $\Delta=4$ hours. The causal
encoder reads the full event stream, actions included. The state of window $w$
is the encoder output at the last observation token in the window,
$s_w = h_{\ell(w)}$; the target state is the corresponding output of the EMA
encoder after the parameter-free LayerNorm, $\tilde s_w$. On full MIMIC-IV,
94\% of crops contain at least one consecutive window pair at $\Delta=4$\,h,
with about 11.5 events per window (24-hour windows left too few pairs per crop).

\paragraph{Actions.}
For a horizon of $k$ windows, the action summary is the mean content
embedding of the action events in windows $w,\dots,w+k-1$, together with the
log action count:
\begin{equation}
a_{w,k} = \Big[\tfrac{1}{|A_{w,k}|}\textstyle\sum_{j\in A_{w,k}} E_{c_j}\;;\;\log(1+|A_{w,k}|)\Big].
\end{equation}

\paragraph{Prediction and loss.}
An MLP predictor $p_\psi$ maps the current state, the actions, the horizon and
the elapsed time to the future state:
\begin{equation}
\hat s_{w+k} = p_\psi\big(s_w,\,a_{w,k},\,e(k),\,e(t_{w+k}-t_w)\big),
\qquad
\mathcal{L}_{\mathrm{state}} = \frac{1}{|\mathcal{K}|}\sum_{k\in\mathcal{K}}
\frac{1}{|\mathcal{P}_k|}\sum_{w\in\mathcal{P}_k}\ell_\beta\big(\hat s_{w+k},\,\tilde s_{w+k}\big),
\end{equation}
with $\mathcal{K}=\{1,2,4\}$ windows (4, 8 and 16 hours ahead) and
$\mathcal{P}_k$ the window pairs present in the crop. The total loss keeps the
hybrid's weights, replacing only the event-level latent target with the
state-level one:
\begin{equation}
\mathcal{L}_{\mathrm{SAC}} = \mathcal{L}_{\mathrm{state}} + 0.1\,\mathcal{L}_{\mathrm{code}} + 0.05\,\mathrm{SIGReg}.
\end{equation}

\paragraph{Controls.}
The \emph{treatment-blind} control is identical except that $a_{w,k}$ is set to
zero, so any difference between the two isolates the use of the action input.
After training we also measure action sensitivity on held-out crops: the
state-prediction error with the true actions versus with actions permuted
across patients. A model that ignores actions shows no difference.

%==============================================================================
\section{Evaluation protocol}
\label{sec:eval}

\paragraph{Tasks.}
Seven binary tasks, defined declaratively in ACES \citep{aces}: 365-day
mortality, 365-day inpatient admission, 30-day readmission (anchored at
discharge), and first-ever diagnosis within 365 days of chronic kidney
disease, COPD, diabetes and heart failure (restricted to patients with no prior
occurrence). One anchor per subject is drawn by a seeded hash among candidate
times with at least 32 prior events, before death, and with complete follow-up
or death within the horizon. The headline metric is the mean AUROC over the six
non-mortality tasks (``mean-of-6''); mortality is reported separately.

\paragraph{Leakage control.}
Model inputs are cut strictly before the anchor at a single choke point; a test
inserts a death one day after each anchor and asserts that features and
embeddings are bit-identical.

\paragraph{Baselines.}
Count features (log counts over 30 days, 365 days and all history, most recent
value per code, numeric-observation counts) feed L2 logistic regression (LR) and
XGBoost (GBM) \citep{xgboost}, both tuned on the tuning split. We audited GBM
for leakage through near-synonyms of outcome codes that the task predicates do not
exclude (e.g.\ hypertensive CKD for CKD, insulin use for diabetes): these appear
in 0.7--2.0\% of positives, carry no measurable share of GBM's gain, and removing
them leaves held-out AUROC unchanged. GBM's most important features are many
weak signals, each at most 0.75\% of total gain.

\paragraph{Frozen probe and fine-tuning.}
The \emph{probe} fits L2 logistic regression on the final valid token's
representation (causal) or the mean over tokens (bidirectional), with $C$ tuned
on the tuning split. \emph{Fine-tuning} trains the encoder with a
LayerNorm\,+\,linear head per task under binary cross-entropy, AdamW
(learning rate $2\times10^{-5}$ for the encoder, $10^{-3}$ for the head), at most five epochs,
early stopping on tuning AUROC with patience 2. A randomly initialized encoder of the
same architecture is the untrained control for both.

\paragraph{Statistics.}
All models are scored on identical anchors (1{,}346 held-out subjects have an
anchor; 1{,}025--1{,}120 rows per task). Differences use a paired bootstrap:
2{,}000 resamples of subjects, applied jointly to both models and all tasks,
95\% percentile intervals. When three seeds exist, the per-resample difference
is averaged over seeds; this does not shrink test-set sampling variance, since
all seeds share the same subjects. The seed-to-seed spread of mean-of-6 for one
configuration at 300M tokens is 0.009; we do not interpret differences below it
without a paired interval.

\paragraph{Compute.}
Every model is pretrained for 1B nominal tokens (30{,}518 steps of $64\times512$;
padding makes this about 550M real events) unless stated, on one RTX~4060
(8\,GB): about 2.7 hours for AR and 3.6 hours for the hybrid.

%==============================================================================
\section{Results}
\label{sec:results}

\subsection{Latent-only objectives lose to a token loss}

\begin{table}[t]
\centering\small
\caption{Objective families on the 8\% pool at 300M tokens, frozen probe,
mean-of-6 AUROC, mean of three seeds (seed values in parentheses).}
\label{tab:objectives}
\begin{tabular}{lll}
\toprule
objective & mean-of-6 & seeds \\
\midrule
hybrid (next-latent + 0.1 next-code)      & \textbf{0.770} & (0.772, 0.766, 0.773) \\
next-code AR                              & 0.768 & (0.775, 0.764, 0.765) \\
masked-span JEPA, EMA target              & 0.705 & (0.713, 0.705, 0.697) \\
masked-span, code reconstruction only     & 0.700 & (0.700, 0.701, 0.700) \\
window-pooled latent (30/365 d)           & 0.700 & (0.701, 0.701, 0.700) \\
next-latent, no code loss                 & 0.682 & (0.675, 0.688, 0.682) \\
\midrule
untrained encoder (control)               & 0.666 & --- \\
\bottomrule
\end{tabular}
\end{table}

Table~\ref{tab:objectives} compares six objectives at a 300M-token budget.
Every objective without a dense next-code loss on the encoder trails AR by
0.06--0.09, far outside the 0.009 seed spread. Latent prediction is not empty:
masked-span JEPA beats the untrained control by $+0.039$. But the gap to a
token loss is large, consistent with the MIMIC-IV finding of
\citet{imperialjepa} that JEPA trails MLM. At this budget the hybrid and AR are
tied.

\subsection{The hybrid separates from AR with more tokens}

At 1B tokens on the 8\% pool, probed, AR reaches 0.769 (seeds 0.772, 0.770,
0.767) and the hybrid 0.794 (0.798, 0.791, 0.794): a $+0.025$ gap with
non-overlapping seed ranges. At 2B tokens neither improves (hybrid 0.798, AR
0.764, seed 0); at 1B tokens the 8\% pool is already seen about 12 times, so we
scale patients rather than tokens from here on. LR and GBM reach 0.795 and
0.808 on the same split, so the frozen hybrid remains below GBM.

\subsection{Probe versus fine-tuning}

\begin{table}[t]
\centering\small
\caption{Mean-of-6 AUROC on the 8\% held-out subjects with 8\% train labels,
as the pretraining pool grows. 3 seeds where marked (seed range in
brackets); otherwise seed 0. Count baselines do not depend on pretraining.}
\label{tab:main}
\begin{tabular}{llcccc}
\toprule
pretraining pool & objective & frozen probe & fine-tuned & seeds \\
\midrule
\multirow{2}{*}{8\% (24k pts)}  & AR     & 0.769 & 0.806 & 3 \\
                                & hybrid & 0.794 & 0.813 & 3 \\
\addlinespace
\multirow{2}{*}{25\% (73k pts)} & AR     & 0.785$^\dagger$ & 0.813 & 3 (FT) \\
                                & hybrid & 0.797$^\dagger$ & 0.821 \,[0.818, 0.825] & 3 (FT) \\
\addlinespace
\multirow{2}{*}{full (249k pts)} & AR    & 0.788 & 0.814 & 1 \\
                                 & hybrid & 0.803 & \textbf{0.822} & 1 \\
\midrule
none & from scratch (supervised) & --- & 0.689 & 1 \\
\multicolumn{2}{l}{count features: GBM / LR} & \multicolumn{2}{c}{0.808 / 0.795} & --- \\
\bottomrule
\multicolumn{5}{l}{\footnotesize $^\dagger$seed 0.}
\end{tabular}
\end{table}

Table~\ref{tab:main} is the central result. Under the frozen probe, no
encoder at any pool size reaches GBM: the best is the full-pool hybrid at 0.803
against 0.808. Under fine-tuning, every hybrid reaches or passes GBM, and
the ordering hybrid $>$ AR holds at every pool size. Two conclusions drawn from the
probe alone would have been wrong. First, a probe-only study would conclude
that pretrained encoders do not beat a tuned count model. Second, under the
probe, growing the pool from 8\% to 25\% helped AR ($+0.016$) but not the
hybrid ($+0.003$, inside seed spread); under fine-tuning it helped the hybrid
($0.813\to0.821$, above GBM on every seed). The 256-dimensional frozen
readout, not the encoder, was the bottleneck: LR on raw counts (0.795 on 8\% labels,
0.830 when refit on 25\% labels) beats the frozen hybrid probe in both cases.
Fine-tuned gains are flat from 25\% to the full pool at the fixed 1B budget
(about 1.8 passes over the full pool), so the token budget, not the patient
count, is the likely limit there.

\begin{table}[t]
\centering\small
\caption{Per-task held-out AUROC, full-pool pretraining, fine-tuned, seed 0.}
\label{tab:pertask}
\setlength{\tabcolsep}{4pt}
\begin{tabular}{lcccccccc}
\toprule
model & inpat.\ & readm.\ & CKD & COPD & diab.\ & HF & mean-of-6 & mort.\ \\
\midrule
hybrid (FT) & \textbf{0.732} & \textbf{0.631} & \textbf{0.939} & \textbf{0.801} & \textbf{0.925} & \textbf{0.904} & \textbf{0.822} & 0.944 \\
AR (FT)     & 0.717 & 0.628 & 0.937 & 0.796 & 0.914 & 0.893 & 0.814 & \textbf{0.948} \\
GBM         & 0.700 & 0.617 & 0.930 & 0.789 & 0.913 & 0.900 & 0.808 & 0.935 \\
LR          & 0.680 & 0.626 & 0.890 & 0.796 & 0.906 & 0.875 & 0.795 & 0.939 \\
\bottomrule
\end{tabular}
\end{table}

\subsection{What is and is not significant}

\begin{table}[t]
\centering\small
\caption{Paired bootstrap of fine-tuned mean-of-6 differences (2{,}000
resamples over the 1{,}346 held-out subjects; 95\% CI; $P$ = fraction of
resamples $>0$). ``Pooled'' averages the three seeds within each resample.}
\label{tab:boot}
\begin{tabular}{llcc}
\toprule
comparison & pretraining pool & difference [95\% CI] & $P(>0)$ \\
\midrule
hybrid $-$ AR  & 25\%, pooled 3 seeds & $+0.008$ \ci{+0.002}{+0.014} & 0.995 \\
hybrid $-$ AR  & 8\%, pooled 3 seeds  & $+0.007$ \ci{-0.000}{+0.013} & 0.97 \\
hybrid $-$ AR  & full, seed 0         & $+0.008$ \ci{+0.000}{+0.016} & 0.98 \\
\addlinespace
hybrid $-$ GBM & 25\%, pooled 3 seeds & $+0.013$ \ci{-0.001}{+0.028} & 0.97 \\
hybrid $-$ GBM & full, seed 0         & $+0.014$ \ci{-0.001}{+0.028} & 0.97 \\
AR $-$ GBM     & 25\%, pooled 3 seeds & $+0.005$ \ci{-0.009}{+0.020} & 0.76 \\
\bottomrule
\end{tabular}
\end{table}

Table~\ref{tab:boot} gives the paired intervals. Hybrid $>$ AR is
significant with 25\%-pool pretraining and borderline at the other two pool sizes.
Hybrid $>$ GBM is not significant: the interval includes zero in every
configuration, although $P(>0)\ge 0.97$. AR is indistinguishable from GBM. The
limit is the test set: 1{,}346 labelled subjects give intervals about
$\pm0.015$ wide. Scoring the full-pool held-out split ($\approx$36k subjects) is
the direct remedy (\S\ref{sec:limits}).

\subsection{Action-conditioned patient-state gate}
\label{sec:gate}

\begin{table}[t]
\centering\small
\caption{State-model gate: full-pool pretraining, 1B tokens, scored on the 8\%
held-out subjects. Seed 0 / seed 1. \textcolor{red}{Placeholders: experiment in
progress.}}
\label{tab:gate}
\begin{tabular}{lccc}
\toprule
objective & frozen probe & fine-tuned & mortality (FT) \\
\midrule
hybrid                         & 0.803 / \todo{s1} & 0.822 / \todo{s1} & 0.944 / \todo{s1} \\
AR                             & 0.788 / \todo{s1} & 0.814 / \todo{s1} & 0.948 / \todo{s1} \\
MLM (bidirectional, 15\%)      & \todo{} & \todo{} & \todo{} \\
state, treatment-blind         & \todo{} & \todo{} & \todo{} \\
state, action-conditioned      & \todo{} & \todo{} & \todo{} \\
\midrule
GBM / LR                       & \multicolumn{2}{c}{0.808 / 0.795} & 0.935 / 0.939 \\
\bottomrule
\end{tabular}
\end{table}

Table~\ref{tab:gate} reports the gate, run under the same protocol as the fine-tuned results above. The pre-stated pass
condition is that the action-conditioned state model beats the hybrid, AR, MLM
and GBM under fine-tuning, and beats its treatment-blind twin.
\todo{paired bootstrap: state\_ac $-$ hybrid, state\_ac $-$ state\_blind, state\_ac $-$ GBM, MLM $-$ AR}
\todo{action sensitivity: held-out state-prediction error, true vs permuted actions, for state\_ac and state\_blind}
\todo{pass/fail statement against the pre-stated condition}

%==============================================================================
\section{Discussion}
\label{sec:discussion}

\paragraph{Why the probe misleads.}
A frozen linear probe on one pooled vector measures how linearly accessible a
task is in the final representation, not how useful the encoder is as an
initialization. The two can disagree in either direction, and here they
disagree in the direction that matters: the probe said more patients helped AR but not the
hybrid, and that no encoder beat a count model; fine-tuning says the opposite on
both counts. All prior EHR-JEPA comparisons we know of
\citep{smbstructure,clinjepa,imperialjepa} are probe-only. Their rankings may be right, but
the evidence here shows that a probe-only ranking can change under
fine-tuning, so we recommend reporting both.

\paragraph{Why the code loss is needed.}
Latent targets from an EMA encoder are only as informative as that encoder. Without
a token loss, nothing anchors the target to the identity of clinical events,
and the predictor can succeed by predicting the slowly varying, easy part of
the representation: the pure next-latent objective reaches a very low
prediction loss with a near-zero gap between matched and shuffled target similarity, the
signature of copying. The code loss keeps the representation discriminative; the latent
term then adds a multi-horizon prediction signal the next-code loss does not
provide, and SMB-Structure's improvements under a hybrid are consistent with this account.
Our results do not support the stronger claim that latent prediction alone is
a better EHR objective.

\paragraph{From events to states and actions.}
Event-level latent targets mix what happens to the patient with what clinicians
decide to do \citep{imperialjepa}. Conditioning on treatments separates the two
in the predictor's inputs. The resulting model is a predictive model of
future states given observed treatments; because treatments in observational
data are chosen in response to the patient's state, its action-conditioned
predictions are confounded and should not be read as treatment effects. Whether this helps downstream
prediction is what the gate tests. \todo{interpretation once gate results are in}

\paragraph{Planned: a trial-anchored benchmark.}
A natural external check for an action-conditioned model is whether its
predictions, used inside a target-trial emulation \citep{hernan2016}, reproduce
effects established by randomized trials \citep{rctduplicate}; this requires
adjustment for confounding by indication, which our model does not provide. We plan a public
benchmark on MIMIC-IV, and on EHRSHOT and INSPECT subject to access, that
pairs standard prediction tasks with such trial-anchored
checks. To our knowledge no public benchmark yet combines structured EHRs,
action-conditioned models and trial-anchored checks; this is future work, and no
results are reported for it here.

%==============================================================================
\section{Limitations}
\label{sec:limits}

\begin{itemize}[leftmargin=*,itemsep=2pt,topsep=2pt]
\item \textbf{Test-set size.} All downstream numbers use 1{,}346 labelled
held-out subjects; this is why hybrid $-$ GBM is not significant. Scoring
on the full-pool held-out split is the next step.
\item \textbf{Seeds.} Full-pool results (Table~\ref{tab:pertask}) are one seed;
the 25\%-pool fine-tuned results are three.
\item \textbf{One site, one model size.} MIMIC-IV is a single academic medical
centre with ICU and hospital data. All models have 10M parameters; a 35M run was
started and stopped before completion. EHRSHOT and INSPECT are planned.
\item \textbf{Fixed token budget.} Full-pool models see each event about 1.8
times; fine-tuned gains are flat from 25\% to the full pool at that budget.
\item \textbf{Baselines.} We compare with count-based GBM/LR and our own
reimplementations of AR and MLM, not with released pretrained models such as
SMB-Structure, CLMBR or MOTOR, which differ in data and size.
\item \textbf{Actions are code-defined and confounded.} The action set is a
code-list heuristic; dose, route and timing within a window are not modelled.
Treatments are assigned in response to patient state, so the state model is
predictive only; it supports no causal or counterfactual claim.
\end{itemize}

%==============================================================================
\section{Conclusion}

On public MIMIC-IV data, with seeds, paired intervals and a leakage-checked
tabular baseline, latent prediction does not replace a token loss for EHR
pretraining but improves on it as an auxiliary term, and the improvement is
visible under fine-tuning and partly hidden under the frozen probe that
previous EHR-JEPA studies relied on. The margin over tuned gradient boosting is
positive but not significant at our test-set size. We tested a small code-token
action-conditioned patient-state objective, with a treatment-blind control, to see whether
the patient-state formulation proposed for EHR JEPAs, with treatments as inputs,
improves on event-level targets. \todo{gate conclusion}

\paragraph{Reproducibility.}
Code (Apache-2.0), grid configurations and per-run summaries are in the
project repository; every table cell maps to a committed configuration.
\todo{repository URL and commit hashes}

\bibliographystyle{plainnat}
\bibliography{refs}

\end{document}
